% Compilar con XeLaTeX o LuaLaTeX.
\documentclass[11pt,letterpaper]{article}

\usepackage{iftex}
\ifPDFTeX
  \usepackage[utf8]{inputenc}
  \usepackage[T1]{fontenc}
  \usepackage{lmodern}
  \usepackage[spanish]{babel}
\else
  \usepackage{fontspec}
  \usepackage{polyglossia}
  \setmainlanguage{spanish}
\fi
\usepackage[margin=2.3cm]{geometry}
\usepackage{amsmath,amssymb}
\usepackage{enumitem}
\setlength{\emergencystretch}{3em}
\pagestyle{plain}

\newcounter{preg}
\newcommand{\pregunta}[1]{\refstepcounter{preg}\par\medskip\noindent\textbf{P\thepreg.\ #1}\par\nopagebreak\smallskip}
\newcommand{\RJ}{R_{\mathrm J}}

\title{Reto III.3: banco de preguntas para la defensa}
\author{Frankly Giovanni Poveda Pinzón \and Camila Castillo López}
\date{noviembre de 2026}

\begin{document}
\maketitle

\noindent Preguntas probables sobre la física, el código, los datos y las decisiones, con una
respuesta corta. Las marcadas con (T) conviene saber hacerlas en el tablero.

\section{Formulación y física}

\pregunta{¿Por qué la sonda gana energía si la gravedad es conservativa?}
Porque la energía que se conserva es la del sistema completo. En el sistema de Júpiter la
sonda sale con la misma rapidez $v_\infty$ con la que llegó; solo cambia la dirección. Visto
desde el Sol, $\vec v = \vec v_{\mathrm J} + \vec v_\infty$, y al girar $\vec v_\infty$ cambia
$|\vec v|$. En el sistema del Sol el potencial de Júpiter depende del tiempo (Júpiter se mueve),
así que la energía de la sonda no tiene por qué conservarse. La energía la cede Júpiter, cuyo
cambio es imperceptible por la razón de masas.

\pregunta{(T) Escriba el lagrangiano en el sistema rotante y explique cada término.}
$L = \tfrac12[(\dot x - ny)^2 + (\dot y + nx)^2 + \dot z^2] + GM_1/r_1 + GM_2/r_2$. Los
paréntesis son la velocidad inercial $\dot{\vec r} + \vec\omega\times\vec r$ con
$\vec\omega = n\hat z$. Desarrollado: $T_2 = \tfrac12\dot r^2$ (cinética relativa),
$T_1 = n(x\dot y - y\dot x)$ (da Coriolis) y $T_0 = \tfrac12 n^2(x^2+y^2)$ (centrífugo).

\pregunta{(T) Deduzca la ecuación de movimiento en $x$.}
$\partial L/\partial\dot x = \dot x - ny$, su derivada temporal $\ddot x - n\dot y$;
$\partial L/\partial x = n(\dot y + nx) + \partial_x(GM_1/r_1 + GM_2/r_2) = n\dot y + \partial_x\Omega$.
Restando: $\ddot x - 2n\dot y = \partial_x\Omega$.

\pregunta{¿Qué se conserva en el sistema rotante y por qué?}
La función de energía $h = \sum\dot q\,\partial L/\partial\dot q - L$, porque
$dh/dt = -\partial L/\partial t$ y en ese sistema el Sol y Júpiter están quietos, así que $L$ no
depende explícitamente de $t$. Resulta $h = \tfrac12\dot r^2 - \Omega$; $C_J = -2h$ es la constante
de Jacobi.

\pregunta{(T) ¿Por qué no aparece el término de Coriolis en $h$?}
Por el teorema de Euler para funciones homogéneas: $\sum\dot q\,\partial T_k/\partial\dot q = kT_k$.
Entonces $\sum\dot q\,\partial L/\partial\dot q = 2T_2 + T_1$ y $h = 2T_2 + T_1 - (T_2+T_1+T_0-V)
= T_2 - T_0 + V$. Físicamente, la fuerza de Coriolis es perpendicular a la velocidad y no hace
trabajo.

\pregunta{(T) Muestre que $h = E - n\ell_z$.}
$E = T + V = T_2 + T_1 + T_0 + V$ y $n\ell_z = n[x(\dot y + nx) - y(\dot x - ny)] = T_1 + 2T_0$.
Restando: $E - n\ell_z = T_2 - T_0 + V = h$.

\pregunta{Si $h$ se conserva, ¿por qué no se conservan $E$ ni $\ell_z$?}
En el sistema inercial el potencial depende de $\varphi$ y de $t$ solo como $\varphi - nt$. Entonces
$dE/dt = \partial V/\partial t \ne 0$ (no hay simetría temporal) y
$d\ell_z/dt = -\partial V/\partial\varphi \ne 0$ (no hay simetría rotacional). Como
$\partial V/\partial t = -n\,\partial V/\partial\varphi$, $dE/dt = n\,d\ell_z/dt$. La simetría
combinada (trasladar $t$ y rotar $n\Delta t$) da, por Noether, $E - n\ell_z$.

\pregunta{¿Qué es el parámetro de Tisserand y por qué no se conserva en los datos?}
Es $h$ escrito con los elementos heliocéntricos de la sonda, suponiendo Júpiter en órbita circular,
$\alpha\to0$ y la sonda lejos de Júpiter: $T_{\mathrm J} = a/a_s + 2\cos i\sqrt{(a_s/a)(1-e^2)}$.
Con Voyager~2, prolongando la trayectoria hasta más de $3\,r_{\mathrm{SOI}}$ a cada lado, pasa de
2.576 a 2.670 ($3.6\,\%$). No es por estar cerca de Júpiter, sino porque la órbita de Júpiter es
excéntrica: Tisserand se deduce de $h$ con el $n$ medio.

\pregunta{(T) Si la órbita de Júpiter no es circular, ¿qué queda de $\Delta E = n\Delta\ell_z$?}
En el sistema inercial, $dE/dt = \partial V/\partial t = \sum_i\vec a_i\cdot\vec V_i$: la energía
cambia por lo que cada cuerpo acelera a la sonda en la dirección en que ese cuerpo se mueve.
Escribiendo $\vec V_i = \omega\,\hat z\times\vec R_i + \dot R_i\hat R_i$ queda
$dE/dt = \omega(t)\,d\ell_z/dt + \sum_i\dot R_i(\vec a_i\cdot\hat R_i)$. Con órbita circular
($\omega = n$, $\dot R_i = 0$) se recupera Jacobi. Con órbita excéntrica hay que usar la velocidad
angular instantánea $\omega = n\sqrt{1-e^2}\,(a/R)^2$, que para Júpiter va de $0.91\,n$ a $1.10\,n$.

\pregunta{¿Se cumple $\Delta E = n\Delta\ell_z$ con los datos reales de Voyager~2?}
Con el $n$ medio, no: $\Delta E = 164.2$ y $n\Delta\ell_z = 173.2$~km$^2$/s$^2$ ($5.5\,\%$). En julio
de 1979 Júpiter estaba a 5.33~UA, cerca del afelio, y giraba a $0.952\,n$. Con la velocidad angular
instantánea, $\int\omega\,d\ell_z = 164.8$ ($0.4\,\%$). Ese resto se reparte entre la velocidad
radial de Júpiter ($-0.15$) y las fuerzas que no están en el potencial, satélites y maniobra
($-0.48$); el Modelo~2 ampliado da para estas últimas $-0.47$ por otro camino.

\pregunta{¿La sonda quedó en una órbita de escape?}
Sí. Lejos de Júpiter (más de $3\,r_{\mathrm{SOI}}$) la energía heliocéntrica pasa de $-121.8$ a
$+26.3$~km$^2$/s$^2$ y la excentricidad de 0.73 a 1.30. Ojo: al final del registro, todavía dentro
de la esfera de influencia, la energía calculada solo con el Sol es $+51$; ese valor no es el
asintótico porque falta el potencial de Júpiter.

\pregunta{(T) Deduzca el ángulo de giro $\delta$.}
Hipérbola $\rho = (\ell^2/\mu)/(1+e\cos f)$. Las asíntotas: $\cos f_\infty = -1/e$. El ángulo entre
las velocidades de entrada y salida es $\delta = 2f_\infty - \pi$, así que
$\sin(\delta/2) = -\cos f_\infty = 1/e$. Con $\ell = r_pv_p$ y $v_p^2 = v_\infty^2 + 2\mu/r_p$ se
obtiene $e = 1 + r_pv_\infty^2/\mu$.

\pregunta{¿Cómo se gana energía y cómo se pierde?}
$\Delta E = \vec v_{\mathrm J}\cdot\Delta\vec v_\infty$. Si $\vec v_\infty$ gira hacia la dirección de
movimiento de Júpiter (la sonda pasa por detrás del planeta), gana energía; si gira en contra
(pasa por delante), la pierde. Es coherente con $\Delta E = n\Delta\ell_z$, porque
$n\Delta\ell_z = na\,\Delta v_t = v_{\mathrm J}\Delta v_t$.

\pregunta{¿Qué es la esfera de influencia y de dónde sale $a(M_2/M_1)^{2/5}$?}
Laplace compara el cociente perturbación/atracción en los dos sistemas. Desde Júpiter, la marea
solar $\sim GM_1\rho/a^3$ frente a $GM_2/\rho^2$; desde el Sol, $GM_2/\rho^2$ frente a $GM_1/a^2$.
Igualando los cocientes: $\rho^5 = (M_2/M_1)^2a^5$. Para Júpiter son $674\,\RJ \approx 48$
millones de km.

\pregunta{Si en el borde de la esfera el Sol atrae 4 veces más que Júpiter, ¿por qué se ignora?}
Porque en el sistema de Júpiter lo que entra es la \emph{diferencia} entre lo que el Sol acelera a
la sonda y a Júpiter (el término $\vec d/d^3$ cancela la parte uniforme). Esa diferencia es la
marea, de orden $GM_1\rho/a^3$, mucho menor que la atracción total.

\section{Modelo 1 y código}

\pregunta{¿Qué es el problema de Lambert y cómo lo resolvió?}
Dadas dos posiciones y el tiempo de vuelo, hallar la órbita que las une. Se usa la variable
universal $z$ (Curtis, alg.~5.2): $F(z) = 0$ es monótona; se busca un intervalo con cambio de signo
sobre una malla y se refina con el método de Brent. Se verifica propagando la solución con un
integrador numérico (llega al punto pedido) y con el límite de Hohmann.

\pregunta{¿Cómo decide si un sobrevuelo es factible?}
Con los dos tramos de Lambert se obtienen $\vec v_{\infty,\mathrm{ent}}$ y el
$\vec v_{\infty,\mathrm{sal}}$ que exige el segundo tramo. Es factible sin motor si sus magnitudes
coinciden (tolerancia 0.3~km/s en la rejilla y cero en el refinado) y si el $r_p$ que da
$\sin(\delta/2) = 1/e$ es mayor que $2\,\RJ$.

\pregunta{¿Qué es $C_3$ y por qué probó varios valores?}
$C_3 = v_\infty^2$ al salir de la Tierra, la energía que debe dar el lanzador. No conocemos el valor
exacto del cohete de 1977, así que es una hipótesis (H8) y estudiamos la sensibilidad. Con
efemérides reales, las fechas de Voyager~2 requieren $C_3 = 103$~km$^2$/s$^2$.

\pregunta{¿Por qué las órbitas circulares daban un error tan grande?}
Júpiter tiene excentricidad 0.048: su distancia al Sol varía $\pm5\,\%$ y su velocidad también. Eso
cambia $\vec v_\infty$ y el giro necesario. En las fechas de Voyager~2, el $r_p$ requerido pasa de
12.6 (circular) a $9.72\,\RJ$ (DE421); el real fue $10.09\,\RJ$.

\section{Modelo 2 y control numérico}

\pregunta{¿Por qué REBOUND con IAS15?}
Es un integrador de orden 15 con paso adaptativo: da pasos de horas cerca del periapsis y de días
en el crucero, y mantiene el error en la precisión de máquina. Lo controlamos con un integrador
independiente (DOP853 de SciPy).

\pregunta{¿Cómo sabe que la simulación es correcta?}
(1) La constante de Jacobi se conserva a $2.5\times10^{-14}$ mientras la energía heliocéntrica
cambia más de 100~km$^2$/s$^2$. (2) Al cambiar la tolerancia de $10^{-4}$ a $10^{-11}$ la posición
final a 250 días cambia como mucho 2.1~m (y menos de 2~cm con $\epsilon \le 10^{-6}$); DOP853 difiere 3~m. (3) Límites: sin masa de Júpiter no hay
cambio de energía, y cerca del periapsis los dos modelos coinciden.

\pregunta{¿Cómo midió el error de las cónicas? ¿Por qué esa medida?}
Ambos modelos parten del mismo estado en el periapsis y se propagan hacia atrás y adelante hasta
más de $3\,r_{\mathrm{SOI}}$. Se compara la ganancia de energía heliocéntrica $\Delta E$ entre los
extremos. Es la cantidad física que produce el sobrevuelo y, a diferencia de un error de posición a
tiempo fijo, no depende de cuánto se propague después.

\pregunta{¿Por qué el error crece al bajar $v_\infty$?}
Con $v_\infty$ pequeña la sonda tarda más en cruzar la región cercana al borde de la esfera, donde
el Sol y Júpiter tienen efectos comparables; las cónicas ignoran uno u otro todo ese tiempo. Era la
predicción escrita antes de simular.

\pregunta{¿Por qué el error cambia de signo con radios de empalme grandes?}
Con radios pequeños se ignora a Júpiter durante mucho tiempo fuera de la esfera y el error es
grande y positivo (las cónicas sobreestiman la ganancia). Con radios grandes se ignora al Sol dentro
de la esfera, un error de signo opuesto que domina cuando $v_\infty$ es alta. Los dos efectos se
compensan en parte; Laplace y Hill quedan cerca del punto de compromiso.

\pregunta{¿Por qué decía que el Modelo 2 ``violaba'' Jacobi y luego no?}
La primera medida usaba el Sol fijo en el origen, una aproximación: daba variaciones de $10^{-3}$,
imposibles para un invariante exacto. Con la fórmula baricéntrica correcta la variación es
$2.5\times10^{-14}$. Es uno de los errores de la IA registrados en la bitácora.

\section{Datos reales}

\pregunta{¿De dónde salen los datos y cómo sabe que son correctos?}
Planetas: efeméride JPL DE421. Voyager~2: kernel SPICE \texttt{vgr2\_jup230.bsp} de NAIF, leído con
CSPICE. Satélites galileanos: paquete Astronomy Engine. Comprobación con valores publicados: el
kernel da el máximo acercamiento al centro de Júpiter el 9 jul 1979 a las 22:29:01~UTC a
721\,673~km (publicado: 22:29~UTC, 721\,670~km), y las distancias mínimas a los cuatro satélites
coinciden al $0.2\,\%$ o mejor (Ganímedes: 62\,078 frente a 62\,130~km). Para DE421 se comprueban
la posición de la Tierra en J2000 y las distancias e inclinaciones de Júpiter y Saturno.

\pregunta{¿Usó los datos de Voyager para ajustar algo?}
No. Las fechas de las Voyager no entran en la búsqueda; se evalúan después. La trayectoria SPICE
solo se usa para comparar. Por eso cuentan como ``datos no utilizados en el ajuste''.

\pregunta{¿Por qué el Modelo 2 no seguía a Voyager~2 después del periapsis?}
Con Júpiter puntual el modelo acierta el periapsis con 1300~km de error, pero después se separa
68\,000~km en 11 días y no reproduce la caída de la energía jovicéntrica real ($v_\infty$
equivalente de 7.650 a 7.578~km/s). Faltaban dos cosas. Los satélites galileanos y el achatamiento
de Júpiter explican $-0.029$~km/s. El resto, $-0.033$~km/s, es una maniobra: la sonda encendió
sus motores 76 minutos, unas dos horas después del máximo acercamiento. Con todo incluido el
modelo termina en 7.5780~km/s (real: 7.5777) y el tramo posterior a la maniobra se reproduce con
6~km de error en 11 días.

\pregunta{¿Cómo sabe que es una maniobra y no un error de los datos o del modelo?}
Se calculó la aceleración real (derivada numérica de la velocidad del kernel) menos la
gravitacional del modelo con satélites y $J_2$. Es prácticamente cero ($6\times10^{-4}$~mm/s$^2$)
durante todo el encuentro, salvo un intervalo de 76.6 minutos que empieza 130 minutos después del
máximo acercamiento, con un valor constante de 2.83~mm/s$^2$ ($\Delta v = 13.0$~m/s). La gravedad no
produce un escalón así. Además hay una fuente independiente: el libro de NASA \emph{Voyage to
Jupiter} (SP-439, 1980) cuenta que la corrección de trayectoria fue un encendido de 76 minutos
dos horas después del máximo acercamiento.

\pregunta{Entonces, ¿el sobrevuelo de Voyager~2 fue ``sin motor''?}
Casi. La maniobra fue de 13~m/s, el $0.2\,\%$ de $v_\infty$, así que para el diseño (Modelo~1) la
hipótesis de sobrevuelo balístico es buena: predice el periapsis con $3.7\,\%$ de error. Pero para
seguir la trayectoria real a nivel de kilómetros hay que incluirla.

\section{Decisiones e IA}

\pregunta{¿Cuál fue la hipótesis más importante?}
Las órbitas planetarias: pasar de circulares a reales redujo el error en el periapsis de Voyager~2
de 25\,\% a 3.7\,\%. Para el error de las cónicas, la más importante es el radio de empalme.

\pregunta{¿Qué delegó a la IA y qué no?}
Se delegó lo verificable: código, álgebra intermedia, búsqueda de fuentes de datos y borradores.
Cada resultado se aceptó solo con una verificación independiente. La elección del enfoque, la
medida del error y la interpretación son del equipo. La formulación lagrangiana la redactó la
IA por tiempo; se verificó con sympy y numéricamente, y el equipo la re-derivó.

\pregunta{Dé un ejemplo de error de la IA y cómo lo detectó.}
El solucionador de Lambert fallaba en transferencias de más de \(180^\circ\): una prueba que
propaga la solución con un integrador numérico mostró que no llegaba al punto pedido. Otro: la IA
concluyó que con $C_3\leq 85$ se reproducía a Voyager~2; con efemérides reales se vio que se
necesitan 103~km$^2$/s$^2$. El más instructivo: la IA comprobó $\Delta E = n\Delta\ell_z$ con Voyager~2
y obtuvo un acuerdo del $0.6\,\%$. Era fortuito: omitió el potencial de Júpiter y usó el $n$ medio,
dos errores del $5\,\%$ que se compensaban. Un resultado que sale bien también hay que revisarlo.

\pregunta{¿Quién encontró los errores de la última revisión?}
La auditoría del 10 de noviembre se delegó al asistente, con la instrucción de no dar por bueno
nada de lo anterior. Encontró ocho errores de su propio trabajo. Lo que los hace creíbles no es
quién los encontró sino la evidencia, que no depende del código: horas y distancias publicadas de
los sobrevuelos, un balance de energía que cierra, y una fuente histórica para la maniobra.

\pregunta{Si otro equipo obtiene un error distinto para las cónicas, ¿quién tiene razón?}
Ambos pueden tenerla: el número depende de la medida del error (aquí, $\Delta E$ desde el mismo
periapsis), del radio de empalme y de $r_p$. Lo importante es que la medida esté declarada y
verificada. Nuestro resultado está acotado por la convergencia (metros) y por el invariante de
Jacobi ($10^{-14}$).

\end{document}
